\documentclass[12pt,a4paper]{article}

\usepackage[T2A]{fontenc}
\usepackage[utf8]{inputenc}
\usepackage[russian]{babel}
\usepackage{setspace}
\usepackage{geometry}
\usepackage{enumitem}

\geometry{left=30mm}
\geometry{right=10mm}
\geometry{top=20mm}
\geometry{bottom=20mm}

\onehalfspacing

\begin{document}
	
	\begin{itemize}[label=--]
		\item Сервисные роботы для потребительского использования продемонстрировали устойчивый рост на 11~\%, и в 2024 году было продано около 20{,}1 млн единиц.
		
		\item Снабжение роботов функциональностью ведения диалога на естественном языке с пользователем способствует повышению эффективности человеко-машинного взаимодействия и снижению затрат на обучение пользователей работе с программами и роботами.
		
		\item В области аффективной робототехники наблюдается, с одной стороны, перенос людьми своих ожиданий с межличностной коммуникации на человеко-машинную и, с другой стороны, использование мультимодальной коммуникации, в том числе аффективной.
		
		\item Согласно существующим исследованиям, люди положительно воспринимают эффект персонализации роботов-собеседников, в частности, предпочитают более экстравертных роботов.
		
		\item Исследования в области человеко-машинной коммуникации показывают, что люди при взаимодействии с роботами используют те же социальные и коммуникативные механизмы, что и при взаимодействии друг с другом.
		
		\item Аффективные роботы находят своё применение в медицине, образовании и сфере работы с клиентами.
	\end{itemize}
	
\end{document}
